% page 498 du fac-simile (image p-497 du scan)
% bloc 162.6 x 242.0 mm ; marge gauche 8.0 mm
\pgc{-12.700mm}{0.000mm}{
\l{\cel{3}490}
\l{}
\l{relatar. (netrans. ad).(\sou{pri termini, muziko-soni, afecion}i,}
\l{\cel{3}\sou{interesti)}.\cel{2}Inter-korespondar per adminime kelka traiti}
\l{\cel{3}komuna. - (\sou{relate}\cel{1}: antonimo di \textquotedbl{}\sou{absolute}\textquotedbl{}). DEFIRS.}
\l{}
\l{relayo. I. (a) Loko ube hundi qui vartas, en la parkureyo}
\l{\cel{3}di chasantaro, remplasos la hundi qui esos fatigita ;}
\l{\cel{3}(b) Loko ube kavali quin onu vartigas, ye intervali, en}
\l{\cel{3}posto-voyo, remplasos la kavali qui esos fatigita.-}
\l{\cel{3}II. (\sou{telegrafo)}\cel{2}Aparato di qua la rolo esas lasar la elek-}
\l{\cel{3}tro-fluo pasar, od interruptar ica, kande ula cirkonstan-}
\l{\cel{3}ci esos realigita en altra cirkuito. - - DEFIRS.}
\l{}
\l{\sou{relegar}. (\sou{trans.}) I. (\sou{yuro-cienco, che la Romani antiqu}a).}
\l{\cel{3}Exilar aden loko determinita, sen privaco di la yuri}
\l{\cel{3}civila e politikala. - II. (\sou{nun, en Francia}). Kondampar}
\l{\cel{3}ad enterneso en kolonio. - DEFIS.}
\l{}
\l{\sou{relevar}. (trans.)\cel{2}Determinar la situeso topografiala di}
\l{\cel{3}(ulo) ed indikar lu sur mapo. - FIS.}
\l{}
\l{\sou{reliefo}. La parto qua salias ek la surfaco cirkondanta. DEFIRS}
\l{}
\l{\sou{religio}. Ensemblo ek la kredi e praktiki e ritui, qui havas}
\l{\cel{3}kom skopo homajar deo o dei. - DEFIRS.}
\l{}
\l{reliquio.\cel{2}(religio kristan.) To quo restas de santo pos lua}
\l{\cel{3}morto, ek lua korpo, ek to quon lu esis, ek la instrumenti}
\l{\cel{3}di lua tormenteso. - DEFIS.}
\l{}
\l{\sou{remarkar}.\cel{2}\sou{(trans.}) Atencar ne-vole (ulu, ulo) ne-ordinara,}
\l{\cel{3}ne-kustumala, ne-expektita. - EF.}
\l{}
\l{\sou{remar}. (\sou{netrans}.) Uzar, por igar batelo (ed, olim, navo)avancar,}
\l{\cel{3}longa ligno-stango, kun planketo ye la extremajo per qua}
\l{\cel{3}onu quaze plaudas o batas la aquo e, tale, obtenas la pulso-}
\l{\cel{3}forco. - FIS.}
\l{}
\l{\sou{remediar}. (\sou{trans., per}). Aplikar la kozo qua esas necesa por}
\l{\cel{3}risanigar (fiziologiale o psikale). - (\sou{anke metaf.}) EFIS}
\l{}
\l{\sou{reminicar}. (\sou{trans}.)(\sou{Lor kompozo literaturala, muzikala, artala})}
\l{\cel{3}Memorar fragmenti di ta o ca literaturajo, muzikajo,}
\l{\cel{3}pikturo, statuo, quan onu imitas ne-vole, nekoncie.- F.}
\l{}
\l{\sou{remisar}. I.(\sou{trans.)}\cel{2}Dispensar (ulu) de parto di lua debo, di}
\l{\cel{3}lua puniseso-grado, di peko. - II. (\sou{netrans}.) (\sou{patol}.)}
\l{\cel{3}\sou{(aludante la}\cel{1}simptomi di morbo).\cel{2}Cesar, o diminutar, dum}
\l{\cel{3}kelka tempo.- EFIS.}
\l{}
\l{\sou{remonstrar}. (\sou{netr., pri}). I. Parolar, diskursar, por demons-}
\l{\cel{3}trar (ad ulu) lua nejusteso, la detrimento quan lu agis.}
\l{\cel{3}- II.\cel{1}\sou{(olim, en Francia)}. Protestar (kom parlamento)}
\l{\cel{3}a la rejo kontre to quon ta parlamento judikis kom misuzo}
\l{\cel{3}o tro-uzo di lua autoritato rejala. - DEF.}
}
